%%%%%%%%%%%%%%%%%%%%%%%%%%%%%%%%%%%%%%%%%%%%%%%%%
\section{Data Quality Assurance and Curation}\label{sec:data_quality_and_curation}

The corpus assembled in \Cref{sec:data_sourcing_and_taxonomy} is heterogeneous in reliability. Community-contributed species labels vary in accuracy, Wikimedia Commons images carry no bounding boxes and may depict illustrations, taxidermy or museum specimens, and even trusted sources hold low-quality frames. This section describes the staged filtering pipeline that turns the raw corpus into a curated, quality-scored dataset with verified ground truth, and the correction passes required once downstream analysis exposed what the initial design had missed.

%%%%%%%%%%%%%%%%%%%%%%%%%%%%%%%%%%%%%%%%%%%%%%%%%
\subsection{Staged Quality Filtering Pipeline}\label{sec:staged_filtering_pipeline}

Filtering is source-aware and cost-staged rather than uniform. Cheap checks run first on every image, and expensive checks run only on what survives them, and only where the source's reliability profile calls for it. \Cref{fig:quality-pipeline} shows the pipeline, whose four image-level stages feed the species classifier of \Cref{sec:speciesnet_matching}.

\begin{figure}[H]
\centering
\begin{tikzpicture}[
  node distance=0.5cm and 0.9cm,
  stage/.style={rectangle, draw, rounded corners, fill=gray!10, minimum width=2.5cm, minimum height=1.1cm, align=center, font=\small},
  arr/.style={-{Latex[length=2mm]}, thick}
]
\node[stage] (s0) {Stage 0\\Metadata};
\node[stage, right=of s0] (s1) {Stage 1\\Heuristics};
\node[stage, right=of s1] (s2) {Stage 2\\MegaDetector};
\node[stage, right=of s2] (s3) {Stage 3\\VLM rescue\\(Wikimedia only)};
\node[stage, below=1.1cm of s2] (s5) {SpeciesNet\\classification};
\node[stage, left=of s5] (s6) {SpeciesNet\\filtering \& mapping};
\draw[arr] (s0) -- (s1);
\draw[arr] (s1) -- (s2);
\draw[arr] (s2) -- (s3);
\draw[arr] (s2.south) -- (s5.north);
\draw[arr] (s3.south) |- ($(s3.south)+(0,-0.5)$) -| (s5.north east);
\draw[arr] (s5) -- (s6);
\end{tikzpicture}
\caption[Staged, source-aware quality filtering pipeline]{Staged, source-aware quality filtering pipeline. Cheaper checks run first and gate access to the more expensive stages. The VLM rescue stage applies only to Wikimedia-sourced images that failed automatic animal detection.}
\label{fig:quality-pipeline}
\end{figure}

\paragraph{Stage 0: Metadata}
This stage needs no image I/O and applies per-source criteria to source metadata: a minimum classifier score and detection confidence for GBIF, a research-grade or still-pending community identification for iNaturalist, and a minimum shorter side and photographic file type for Wikimedia. The criteria for the two largest sources were not in fact applied, because a bulk rename of the downloaded files severed the identifier mapping they depend on, which left rejection to the later and more expensive stages.

\paragraph{Stage 1: Heuristics}
Cheap CPU-only checks then apply to every image: corruption, a shorter side of at least $256\,\mathrm{px}$, an aspect ratio no greater than $4.0$, a grayscale test on mean saturation, and a Laplacian blur variance of at least $100$. The $256\,\mathrm{px}$ floor is calibrated against iNaturalist's standard derivative, whose shorter side is roughly $281\,\mathrm{px}$.

\paragraph{Stage 2: MegaDetector}
\textit{MegaDetector} v5 \cite{beeryEfficientPipelineCamera2019}, a general-purpose animal detector, runs on every image surviving Stage 1 and removes content that passes the cheaper heuristics without being a photograph of a live animal, such as drawings, taxidermy mounts and plush toys.

Its output is consumed under two thresholds answering different questions. Admission is strict: an image is retained only if it carries an animal detection above confidence $0.5$ covering at least $1\%$ of the frame. Recording is permissive: every detection above confidence $0.2$ is written to the image's record. Separating the two keeps information a single-threshold design would have discarded, and three later corrections in this section depend on it.

\paragraph{Stage 3: VLM Rescue (Wikimedia Only)}
A rescue stage runs on Wikimedia images that failed Stage 2, the one source holding both unusual real photographs a detector may miss and much illustration and museum material. A caption from \textit{Florence-2-large} \cite{xiaoFlorence2AdvancingUnified2023} is scanned for rejection keywords, and an image containing none is reinstated.

\paragraph{Caption-Based Evaluation}
All surviving images are captioned with \textit{Florence-2-base} \cite{xiaoFlorence2AdvancingUnified2023}, and a language model judges each caption on whether the subject is a real, living animal, whether the image is a photograph rather than a rendering, and whether a substantial part of the animal is visible. The instruction also enumerates what is not grounds for rejection (\Cref{sec:appendix_caption_settings}). A manual review pass over the unsettled cases examined $17{,}737$ images, approving $13{,}827$ and declining $3{,}910$. The $22\%$ rejection rate confirms genuinely ambiguous material, and its per-class availability estimates are what the band structure of \Cref{sec:band_structure} was built on.

%%%%%%%%%%%%%%%%%%%%%%%%%%%%%%%%%%%%%%%%%%%%%%%%%
\subsection{Species-Level Classification and Taxonomic Matching}\label{sec:speciesnet_matching}

Passing the quality gates establishes that an image depicts a real animal, not which of the $225$ classes it shows. \textit{SpeciesNet} \cite{gadotCropNotCrop2024a}, an \textit{EfficientNetV2-M} \cite{tanEfficientNetV2SmallerModels2021} classifier, therefore runs on every MegaDetector-cropped animal region.

Because the taxonomy mixes species-, genus- and family-level labels, a prediction is accepted on a match level read relative to the granularity of the target class rather than absolutely. A genus-level class cannot be matched at species level even by a perfect classifier, so a genus-level prediction there counts as a full match. A coarser prediction counts as partial, and disagreement at family level or above is evidence of a mislabel.

The thresholds are a minimum MegaDetector confidence of $0.5$, a minimum SpeciesNet top-1 score of $0.3$, and a family-mismatch confidence threshold of $0.5$. The third inverts the usual role of confidence: for a family-level disagreement, higher classifier confidence makes rejection more likely, since a confident prediction of a different family is evidence that the source label is wrong.

\begin{table}[H]
\centering
\caption{Outcome of the SpeciesNet-based filter over $465{,}130$ classified records.}
\label{tab:speciesnet-filter-attrition}
\begin{tabular}{@{}lrr@{}}
\toprule
Outcome & Records & Share \\
\midrule
Passed & $158{,}667$ & $34.1\%$ \\
\midrule
No taxonomic match & $164{,}016$ & $35.3\%$ \\
Low classifier confidence & $66{,}769$ & $14.4\%$ \\
Match only at class level & $29{,}969$ & $6.4\%$ \\
Confident family mismatch & $19{,}527$ & $4.2\%$ \\
Match only at order level & $17{,}408$ & $3.7\%$ \\
Predicted taxon outside the $225$ classes & $6{,}218$ & $1.3\%$ \\
Primary crop too small to classify & $2{,}556$ & $0.5\%$ \\
\bottomrule
\end{tabular}
\end{table}

The dominant failure in \Cref{tab:speciesnet-filter-attrition} is no taxonomic match at all rather than a confident disagreement, the first sign that the filter measured the classifier's coverage as much as the data's correctness. \Cref{sec:speciesnet_coverage_gaps} takes that up.

%%%%%%%%%%%%%%%%%%%%%%%%%%%%%%%%%%%%%%%%%%%%%%%%%
\subsection{Correcting SpeciesNet Coverage Gaps}\label{sec:speciesnet_coverage_gaps}

Applying the match-level filter uniformly cost several classes the overwhelming majority of their images, among them \textit{Saguinus} tamarins ($1{,}065 \rightarrow 5$) and pikas ($1{,}484 \rightarrow 31$). Inspecting those cases showed not mislabeling at the source but a coverage gap in \textit{SpeciesNet}, which does not reliably recognize those taxa.

A fallback rule therefore applies where the evidence points that way. If a class's dominant failure is a missing or class-level match and fewer than $20\%$ of its images pass, or an order-level match with fewer than $15\%$ passing, the class bypasses the SpeciesNet filter and falls back to the pre-classification quality-pass flag. This recovered approximately $14{,}500$ images across $17$ classes. Two classes whose dominant failure was instead a confident family mismatch were held under the stricter filter, since that signals label contamination rather than a coverage gap.

%%%%%%%%%%%%%%%%%%%%%%%%%%%%%%%%%%%%%%%%%%%%%%%%%
\subsection{Ground-Truth Annotation Integrity}\label{sec:gt_annotation_integrity}

\paragraph{The Multi-Box Annotation Bug}
A later audit found that every one of $63{,}796$ images in the test annotation file carried exactly one ground-truth bounding box, although MegaDetector had detected multiple animals in a substantial fraction of them. % annotations_test.json
The fix forwards the full detection list through the pipeline and emits one annotation per detection at or above confidence $0.5$, for every image.

\paragraph{Cross-Species Contamination in Multi-Box Images}
Fixing the multi-box bug exposed a subtler defect: every box on an image had been given the image's primary label, so a second animal of a different species was mislabeled as the primary one. The per-box SpeciesNet predictions (\Cref{sec:speciesnet_matching}) were reused to flag secondary boxes automatically, subject to a taxonomic tolerance. Agreement at species, genus or family level counts as acceptable variation. Divergence at order level or above, or no match at all, is flagged instead, provided the secondary box carries a SpeciesNet confidence of at least $0.3$. Without that tolerance the queue would have been unmanageable, since $32{,}401$ of $57{,}205$ multi-detection images carry a differing species index on a secondary box.

The filter flagged $11{,}354$ images, of which $7{,}271$ were resolved: $6{,}477$ as keep, $794$ as edit, and none as discard. No images were removed, and $1{,}344$ mislabeled annotations were deleted across the three splits. The review also yielded $1{,}733$ candidate confusion pairs, which informed the frozen look-alike grouping used for coarse-grained evaluation (\Cref{sec:lookalike_groups}).

%%%%%%%%%%%%%%%%%%%%%%%%%%%%%%%%%%%%%%%%%%%%%%%%%
\subsection{Quality Scoring and Split Assembly}\label{sec:quality_scoring_splits}

Images surviving all stages carry a composite quality score $Q \in [0,1]$ guiding validation sampling and tie-breaking in split assembly:
\begin{equation}
    Q = 0.30\,S_{\text{area}} + 0.25\,S_{\text{edge}} + 0.20\,S_{\text{single}} + 0.25\,S_{\text{conf}}.
    \label{eq:quality-score}
\end{equation}

\begin{table}[H]
\centering
\caption{Components of the quality score of \Cref{eq:quality-score}.}
\label{tab:quality-score-components}
\begin{tabularx}{\textwidth}{@{}lX@{}}
\toprule
Term & Definition \\
\midrule
$S_{\text{area}}$ & $1.0$ at $4$--$40\%$ frame coverage, decaying to $0$ below $2\%$ or above $70\%$ \\
$S_{\text{edge}}$ & Penalizes clipping, reaching $0$ once the box leaves the frame \\
$S_{\text{single}}$ & $1.00$ one detection, $0.60$ two, $0.30$ three or more, $0.50$ none \\
$S_{\text{conf}}$ & MegaDetector box confidence, $0.30$ when unavailable \\
\bottomrule
\end{tabularx}
\end{table}

\Cref{tab:quality-score-components} defines the four components. Two hard exclusions override the score: a box covering less than $1\%$ of the frame, or one extending beyond the image boundary, forces $Q = 0$. An image with no detection at all is dropped rather than kept as a hard example, since it carries no ground-truth box and could only produce false positives.

The three splits are then drawn under deliberately different rules. The test set is drawn first and stratified so that its source mixture mirrors the class's actual pool, since it estimates behavior on the operating distribution. Training images are then taken greedily by quality score from what remains. Validation is sampled from the $30$th to $70$th percentile of $Q$ under a fixed seed, because training has skimmed the top and a random validation set would be easier than the data the model will face.

For the largest classes (Band D, \Cref{sec:band_structure}) a soft source-diversity cap limits any single source to at most $60\%$ of a class's training images, so that one source's photographic style does not become the feature the model learns. The background and human images sit outside this logic, split in fixed proportions. What these rules produce, class by class, is the dataset the rest of the work trains on.

%%%%%%%%%%%%%%%%%%%%%%%%%%%%%%%%%%%%%%%%%%%%%%%%%
\subsection{The Resulting Dataset}\label{sec:final_dataset_composition}

\Cref{tab:final-real-composition} gives the authoritative composition of the curated real-image dataset after the ground-truth corrections of \Cref{sec:gt_annotation_integrity}. Figures reported earlier in this section for intermediate stages are superseded by it.

\begin{table}[H]
\centering
\caption[Composition of the curated real-image dataset by band]{Composition of the curated real-image dataset by band. \enquote{Pool} is the number of quality-passed images available to a class. \enquote{Surplus} is pool imagery beyond the training and test caps, retained but unallocated.}
\label{tab:final-real-composition}
\begin{tabular}{@{}lrrrrrr@{}}
\toprule
Band & Classes & Pool & Train & Val & Test & Surplus \\
\midrule
A & $50$ & $5{,}902$ & $0$ & $0$ & $5{,}803$ & $0$ \\
B & $26$ & $10{,}912$ & $2{,}210$ & $514$ & $8{,}025$ & $0$ \\
C & $26$ & $18{,}517$ & $4{,}420$ & $780$ & $12{,}965$ & $0$ \\
D & $122$ & $422{,}268$ & $137{,}561$ & $11{,}125$ & $36{,}641$ & $231{,}107$ \\
Negatives & $2$ & $2{,}174$ & $1{,}622$ & $126$ & $426$ & $0$ \\
\midrule
Total & $226$ & $459{,}773$ & $145{,}813$ & $12{,}545$ & $63{,}860$ & $231{,}107$ \\
\bottomrule
\end{tabular}
\end{table}

In the exported annotation files this is $145{,}728$ training images with $187{,}705$ boxes, $12{,}543$ validation images with $19{,}732$ boxes, and $63{,}802$ test images with $92{,}094$ boxes, for $222{,}073$ real images and $299{,}531$ annotations across $225$ categories. The real images carry no hand-drawn ground truth. Every box originates from automatic detection, and human effort went into verifying labels rather than drawing boxes.

Band A contributes no training images by design, its classes being supported entirely by synthetic imagery (\Cref{sec:band_structure}). Band D's $231{,}107$ surplus images exceed the per-class caps, and they are what made the enlarged test set of the generator comparison (\Cref{sec:generator_comparison_classes}) possible without touching training data.

\paragraph{Known Defects in the Split}
Three defects in the assembled split are recorded here. First, $99$ hard-excluded images reached the Band A test set, because a Band A class assigns its whole pool to testing without quality filtering. Second, band assignment used a conservative estimate of each class's pool, so some classes sit one band below their actual imagery. Third, the Band B and C test sets are quality-biased low by construction, since training takes the top of the quality distribution and validation the middle, leaving the remainder to test. The effect is bounded but real. It means absolute accuracy on those bands is not directly comparable with Band D's.

\paragraph{A Note on Tuning}
The quality-score weights and breakpoints, the band boundaries, the per-band quotas and the validation window are fixed design choices, justified by reasoning of the kind given above but not selected by any sensitivity analysis. No ablation over them was run. Band A's total dependence on synthetic imagery, visible in \Cref{tab:final-real-composition}, is what \Cref{sec:synthetic_data_supplementation} takes up.
